\begin{figure}

\includegraphics[width=0.85\linewidth]{%
binder-2025-08-24-keyfig/binder/teeplots/2025-08-24-keyfig/baseline=without+bbg=prefatory+subject=population+viz=signchangedotplot+ext=.pdf}%
\hspace{-2.2ex}
\begin{minipage}[t]{0.18\linewidth}%

\vspace{-17.84ex}

\includegraphics[width=\linewidth,trim={0 1.5cm 0 0}, clip]{%
binder-2025-08-24-keyfig/binder/teeplots/2025-08-24-keyfig/baseline=without+bbg=prefatory+subject=population+viz=signchangediffswarmplot+ext=.pdf}

\vspace{-1.5ex}

\includegraphics[width=\linewidth,trim={0 0 0 3.0cm}, clip]{%
binder-2025-08-24-keyfig/binder/teeplots/2025-08-24-keyfig/baseline=without+bbg=prefatory+subject=population+viz=signchangediffswarmplot+ext=.pdf}
\end{minipage}

\vspace{-2ex}
\caption{
\textbf{Fitness trade-offs driven by ecological context.}
\footnotesize
Plots report influence on focal strain competition assays attributable to presence versus absence of background strain.
Experiments competed focal strain samples (population from stint $n+1$) against baseline population from stint $n$.
Outcomes recorded between $1$ (fitness gain; stint $n$ baseline driven extinct) and $-1$ (fitness loss; stint $n+1$ driven extinct), with near-zero values indicating comparable fitness.
Vertical line segments mark discrepancy between abiotic conditions (``without,'' hollow markers) and biotic conditions (filled markers).
Markers omitted where no significant fitness change detected between stint $n+1$ and stint $n$ baseline ($\alpha = 0.005$).
Segment color highlights sign-change effect: harmful trait becomes beneficial under ecological context (blue) or vice versa (red).
Under prefatory context, background strain from stint $n$; under contemporary context, stint $n+1$.
Further detail provided in Supplementary Figure \ref{fig:sign-change}.
}
\label{fig:sign-change}

\end{figure}
